% ============================================================
%  IEEE Conference Paper — California Housing Price Prediction
%  Compiled with: pdflatex → bibtex → pdflatex × 2
%  Required packages: IEEEtran class, standard LaTeX packages
% ============================================================

\documentclass[conference]{IEEEtran}

% ── Core packages ────────────────────────────────────────────
\usepackage[T1]{fontenc}
\usepackage[utf8]{inputenc}
\usepackage{cite}
\usepackage{amsmath,amssymb,amsfonts}
\usepackage{algorithmic}
\usepackage{algorithm}
\usepackage{graphicx}
\usepackage{textcomp}
\usepackage{xcolor}
\usepackage{booktabs}
\usepackage{multirow}
\usepackage{array}
\usepackage{url}
\usepackage{hyperref}
\usepackage{float}
\usepackage{subcaption}
\usepackage{listings}
\usepackage{balance}

% ── Hyperref setup ───────────────────────────────────────────
\hypersetup{
    colorlinks = true,
    linkcolor  = blue,
    citecolor  = blue,
    urlcolor   = blue
}

% ── Code listing style ───────────────────────────────────────
\lstset{
    language     = Python,
    basicstyle   = \ttfamily\footnotesize,
    keywordstyle = \color{blue},
    commentstyle = \color{gray},
    stringstyle  = \color{red!70!black},
    breaklines   = true,
    frame        = single,
    numbers      = left,
    numberstyle  = \tiny\color{gray},
    xleftmargin  = 2em,
    framexleftmargin = 1.5em,
    tabsize      = 4
}

% ── Column separator ─────────────────────────────────────────
\def\BibTeX{{\rm B\kern-.05em{\sc i\kern-.025em b}\kern-.08em T\kern-.1667em\lower.7ex\hbox{E}\kern-.125emX}}

% ============================================================
\begin{document}

\title{%
  Predicting California Median Housing Prices Using\\
  an End-to-End Machine Learning Pipeline:\\
  A Comparative Study of Regression Algorithms\\
  with Automated Feature Engineering and\\
  Hyperparameter Optimisation
}

\author{%
  \IEEEauthorblockN{[Author Name]}
  \IEEEauthorblockA{
    \textit{Department of [Department Name]}\\
    \textit{[Institution Name]}\\
    [City, Country]\\
    [author@email.com]
  }
}

\maketitle

% ============================================================
\begin{abstract}
This paper presents a comprehensive, end-to-end supervised machine learning system for predicting median residential housing prices across California districts using the 1990 United States Census dataset. The pipeline encompasses all canonical stages of applied machine learning: exploratory data analysis, stratified sampling, automated data preprocessing, custom feature engineering, multi-model training and evaluation, and systematic hyperparameter optimisation. Four regression algorithms—Linear Regression, Decision Tree Regressor, Random Forest Regressor, and Gradient Boosting Regressor—are benchmarked under 5-fold cross-validation using Root Mean Square Error (RMSE) as the primary metric. A two-stage hyperparameter search (Randomised Search followed by Grid Search) over the winning Random Forest model yields a final test RMSE substantially lower than the baseline, with a 95\% bootstrap confidence interval computed over 10,000 resamples. Engineered ratio features (\texttt{rooms\_per\_household}, \texttt{bedrooms\_per\_room}, \texttt{population\_per\_household}) are found to improve predictive accuracy, with \texttt{median\_income} emerging as the dominant predictor. The trained pipeline is serialised for deployment and accompanied by a production monitoring stub that issues retraining alerts when live RMSE drifts above a configurable threshold. The complete system demonstrates a reproducible, production-grade approach to regression-based real-estate valuation.
\end{abstract}

\begin{IEEEkeywords}
machine learning, housing price prediction, random forest, gradient boosting, feature engineering, hyperparameter optimisation, California housing dataset, regression, scikit-learn, end-to-end pipeline
\end{IEEEkeywords}

% ============================================================
\section{Introduction}
\label{sec:intro}

Accurate real-estate valuation is a problem of enduring commercial and academic interest. Automated Valuation Models (AVMs) power modern PropTech platforms, inform lending decisions, and guide urban investment strategies. The shift from hedonic pricing models grounded in econometrics \cite{rosen1974hedonic} towards data-driven machine learning (ML) approaches has accelerated since the availability of large, digitised property datasets.

The \emph{California Housing Prices} dataset—derived from the 1990 U.S. decennial census and popularised by Géron \cite{geron2019}—has become a standard benchmark for regression tasks. Its 20,640 district-level records offer a realistic testbed combining geographic, demographic, and housing-stock attributes, together with a continuous target (\texttt{median\_house\_value}) that exhibits non-linear relationships with several predictors.

\subsection{Motivation and Business Context}

The motivating scenario is that of a real-estate investment company that wishes to automate early-stage district screening. A downstream decision system must determine, for each candidate district, whether the predicted median house value warrants further due-diligence investment. The ML model therefore acts as the first-stage filter in a multi-step investment pipeline.

\subsection{Problem Framing}

This is formalised as:
\begin{itemize}
    \item \textbf{Supervised learning}: ground-truth labels (\texttt{median\_house\_value}) are available for every training instance.
    \item \textbf{Multiple univariate regression}: a single continuous target is predicted from nine district-level attributes.
    \item \textbf{Batch inference}: predictions are generated in bulk from stored census snapshots; no streaming or online-learning requirement exists.
\end{itemize}

\subsection{Contributions}

The main contributions of this work are:
\begin{enumerate}
    \item A fully reproducible, open-source end-to-end ML pipeline from raw CSV ingestion to a serialised, deployment-ready artefact.
    \item A custom \texttt{scikit-learn} transformer (\texttt{RatioFeaturesAdder}) that automatically engineers three ratio features shown to improve RMSE.
    \item A systematic comparative evaluation of four regression families under 5-fold cross-validation.
    \item A two-stage hyperparameter optimisation protocol (Randomised Search → Grid Search) that efficiently navigates a high-dimensional configuration space.
    \item A bootstrap confidence interval for the final test RMSE, providing a statistically rigorous estimate of generalisation uncertainty.
    \item A production monitoring stub that enables automated performance degradation detection post-deployment.
\end{enumerate}

The remainder of this paper is structured as follows. Section~\ref{sec:related} reviews related work. Section~\ref{sec:dataset} describes the dataset. Section~\ref{sec:methodology} details the full pipeline. Section~\ref{sec:results} presents experimental results. Section~\ref{sec:discussion} discusses findings and limitations. Section~\ref{sec:deployment} describes deployment considerations. Section~\ref{sec:conclusion} concludes.

% ============================================================
\section{Related Work}
\label{sec:related}

Housing price prediction has been studied through a spectrum of methodological lenses. Early econometric approaches relied on hedonic regression models \cite{rosen1974hedonic}, where property value is decomposed as a function of structural, neighbourhood, and locational attributes. Subsequent work demonstrated that tree-based ensemble methods consistently outperform linear baselines on tabular real-estate data \cite{breiman2001random}.

Pace and Barry \cite{pace1997} applied spatial autoregressive methods to the same California census data, highlighting the importance of geographic structure. More recently, gradient boosting frameworks such as XGBoost \cite{chen2016xgboost} and LightGBM have dominated regression leaderboards on housing datasets, owing to their ability to model complex feature interactions without manual engineering.

Deep learning approaches (fully connected networks, TabNet \cite{arik2021tabnet}) have been applied to property valuation but typically require larger datasets and offer limited interpretability relative to tree ensembles. For the scale and dimensionality of the California housing dataset, Random Forests and Gradient Boosting remain the methods of choice in the literature, consistent with our empirical findings.

The present work aligns most closely with the pedagogical pipeline of Géron \cite{geron2019}, extended with (i) an additional Gradient Boosting competitor, (ii) a two-stage hyperparameter search, (iii) bootstrap confidence intervals, and (iv) a production monitoring framework.

% ============================================================
\section{Dataset Description}
\label{sec:dataset}

\subsection{Source and Provenance}

The California Housing Prices dataset is sourced from the 1990 United States decennial census at the census-block-group (district) level. It was first used by Pace and Barry \cite{pace1997} and subsequently packaged by Géron \cite{geron2019} as a benchmark for regression tasks. The version used here is hosted at \url{https://github.com/ageron/handson-ml2} and downloaded programmatically at runtime.

\subsection{Structure}

\begin{table}[H]
\centering
\caption{Dataset Attributes}
\label{tab:attributes}
\begin{tabular}{@{}llp{3.5cm}@{}}
\toprule
\textbf{Attribute} & \textbf{Type} & \textbf{Description} \\
\midrule
\texttt{longitude}           & Float   & District centroid longitude \\
\texttt{latitude}            & Float   & District centroid latitude \\
\texttt{housing\_median\_age}& Float   & Median age of housing units \\
\texttt{total\_rooms}        & Float   & Total rooms across district \\
\texttt{total\_bedrooms}     & Float   & Total bedrooms (1.0\% missing) \\
\texttt{population}          & Float   & District resident population \\
\texttt{households}          & Float   & Number of households \\
\texttt{median\_income}      & Float   & Median income (scaled, 10k USD) \\
\texttt{ocean\_proximity}    & Categorical & Distance category to ocean \\
\midrule
\texttt{median\_house\_value}& Float   & \textbf{Target} — median house value (USD) \\
\bottomrule
\end{tabular}
\end{table}

Key characteristics:
\begin{itemize}
    \item \textbf{Instances}: 20,640 district records.
    \item \textbf{Missing values}: \texttt{total\_bedrooms} has 207 missing values ($\approx$1.0\%), imputed using column median.
    \item \textbf{Target ceiling}: \texttt{median\_house\_value} is capped at \$500,000 — a known data-collection artefact that inflates density at the upper bound.
    \item \textbf{Categorical levels}: \texttt{ocean\_proximity} takes five values: \texttt{<1H OCEAN}, \texttt{INLAND}, \texttt{NEAR BAY}, \texttt{NEAR OCEAN}, \texttt{ISLAND}.
    \item \textbf{Skewness}: several numerical attributes (\texttt{total\_rooms}, \texttt{total\_bedrooms}, \texttt{population}, \texttt{households}) exhibit strong right-skew, motivating standard scaling.
\end{itemize}

% ── Figure placeholder: Feature Distributions ─────────────────────────────────
\begin{figure}[H]
    \centering
    % Replace with your actual exported figure:
    % \includegraphics[width=\columnwidth]{figures/01_feature_distributions.png}
    \fbox{\parbox{\columnwidth}{\centering\vspace{1.5cm}
    \textbf{[Figure 1 — Feature Distributions]}\\[4pt]
    Insert \texttt{01\_feature\_distributions.png} here.\\
    3×3 grid of histograms (50 bins each) for all nine\\
    numerical features of the training set.\\[4pt]
    Save the figure from your notebook and place it at\\
    \texttt{figures/01\_feature\_distributions.png}
    \vspace{1.5cm}}}
    \caption{Histogram distributions of all numerical features in the training set (16,512 samples, 50 bins). Note the right-skew in count-based attributes and the hard cap at \$500,000 for \texttt{median\_house\_value}.}
    \label{fig:distributions}
\end{figure}

% ============================================================
\section{Methodology}
\label{sec:methodology}

The project follows the eight-step ML project checklist of Géron \cite{geron2019}, implemented entirely in Python~3 using \texttt{scikit-learn}~1.x, \texttt{pandas}, \texttt{numpy}, \texttt{matplotlib}, and \texttt{seaborn}. All randomness is controlled via \texttt{RANDOM\_STATE = 42}.

\subsection{Step 1 — Problem Framing and Performance Measure}

The performance measure is the \emph{Root Mean Square Error} (RMSE):

\begin{equation}
    \mathrm{RMSE}(\mathbf{y}, \hat{\mathbf{y}}) = \sqrt{\frac{1}{m}\sum_{i=1}^{m}\left(y^{(i)} - \hat{y}^{(i)}\right)^{2}}
    \label{eq:rmse}
\end{equation}

where $m$ is the number of instances, $y^{(i)}$ the ground-truth target, and $\hat{y}^{(i)}$ the model prediction. RMSE penalises large errors more heavily than Mean Absolute Error (MAE), making it appropriate for a business context where large under- or over-valuations are costly. Mean Absolute Error (MAE) and the coefficient of determination $R^2$ are reported as supplementary metrics.

\subsection{Step 2 — Data Acquisition and Train/Test Split}
\label{subsec:split}

\subsubsection{Download}
The dataset is fetched from GitHub via \texttt{urllib} and extracted from a \texttt{.tgz} archive into a local directory. Download is idempotent: the script skips re-download if the CSV is already present.

\subsubsection{Stratified Sampling}
Naive random splitting can produce unrepresentative test sets when a key numerical predictor is not uniformly distributed. \texttt{median\_income} is the strongest correlated feature (see Section~\ref{sec:eda}), so it is discretised into five ordinal strata using the bins $[0, 1.5, 3.0, 4.5, 6.0, +\infty)$. \texttt{StratifiedShuffleSplit} (80/20 split) is then applied, ensuring that each income stratum appears in the same proportion in both train and test sets.

\begin{table}[H]
\centering
\caption{Stratification Quality — Income Category Proportions}
\label{tab:stratification}
\begin{tabular}{@{}ccccc@{}}
\toprule
\textbf{Income Cat.} & \textbf{Overall} & \textbf{Train} & \textbf{Test} & \textbf{Train Error (\%)} \\
\midrule
1 & 0.0397 & 0.0397 & 0.0399 & $<$0.1 \\
2 & 0.3188 & 0.3188 & 0.3190 & $<$0.1 \\
3 & 0.3507 & 0.3508 & 0.3501 & $<$0.1 \\
4 & 0.1766 & 0.1766 & 0.1765 & $<$0.1 \\
5 & 0.1142 & 0.1141 & 0.1144 & $<$0.1 \\
\midrule
\textbf{Total} & \textbf{16,512} & \textbf{16,512} & \textbf{4,128} & — \\
\bottomrule
\end{tabular}
\end{table}

After splitting, the temporary \texttt{income\_cat} column is dropped from all DataFrames to prevent data leakage.

\subsection{Step 3 — Exploratory Data Analysis}
\label{sec:eda}

All EDA is performed \emph{exclusively on the training set} to avoid data-snooping bias.

\subsubsection{Geographic Visualisation}

A scatter plot of district coordinates (longitude vs.\ latitude) with marker size proportional to population and colour mapped to \texttt{median\_house\_value} reveals a strong coastal premium: high-value districts cluster along the San Francisco Bay Area, Los Angeles basin, and coastal corridors.

% ── Figure placeholder: Geographic Price Map ──────────────────────────────────
\begin{figure}[H]
    \centering
    % \includegraphics[width=\columnwidth]{figures/02_geo_price_map.png}
    \fbox{\parbox{\columnwidth}{\centering\vspace{1.5cm}
    \textbf{[Figure 2 — Geographic Price Map]}\\[4pt]
    Insert \texttt{02\_geo\_price\_map.png} here.\\
    Scatter plot of (longitude, latitude) coloured by\\
    median house value (RdYlGn colormap), bubble\\
    size proportional to district population.\\[4pt]
    \vspace{1.5cm}}}
    \caption{Geographic distribution of California housing prices. Bubble size is proportional to district population; colour encodes median house value from low (red) to high (green). High-value districts concentrate along coastal regions.}
    \label{fig:geomap}
\end{figure}

\subsubsection{Pearson Correlation Analysis}

Pearson correlation coefficients between each numerical feature and \texttt{median\_house\_value} are computed and ranked:

\begin{table}[H]
\centering
\caption{Pearson Correlation with \texttt{median\_house\_value}}
\label{tab:correlations}
\begin{tabular}{@{}lc@{}}
\toprule
\textbf{Feature} & \textbf{Correlation} \\
\midrule
\texttt{median\_income}      &  $+0.688$ \\
\texttt{total\_rooms}        &  $+0.134$ \\
\texttt{housing\_median\_age}&  $+0.106$ \\
\texttt{households}          &  $+0.065$ \\
\texttt{total\_bedrooms}     &  $+0.049$ \\
\texttt{population}          &  $-0.025$ \\
\texttt{longitude}           &  $-0.045$ \\
\texttt{latitude}            &  $-0.143$ \\
\bottomrule
\end{tabular}
\end{table}

\texttt{median\_income} exhibits by far the strongest positive linear correlation ($r = 0.688$). Count-based attributes (\texttt{total\_rooms}, \texttt{total\_bedrooms}, \texttt{population}, \texttt{households}) have weak linear correlations individually, motivating ratio-based feature engineering.

% ── Figure placeholder: Correlation Bar Chart ─────────────────────────────────
\begin{figure}[H]
    \centering
    % \includegraphics[width=\columnwidth]{figures/03_correlations.png}
    \fbox{\parbox{\columnwidth}{\centering\vspace{1.5cm}
    \textbf{[Figure 3 — Feature Correlation Bar Chart]}\\[4pt]
    Insert \texttt{03\_correlations.png} here.\\
    Horizontal bar chart of Pearson correlation\\
    coefficients, green for positive, red for negative.\\[4pt]
    \vspace{1.5cm}}}
    \caption{Ranked Pearson correlation coefficients between each feature and \texttt{median\_house\_value}. Bars coloured green (positive) and red (negative). \texttt{median\_income} dominates with $r = 0.688$.}
    \label{fig:corr}
\end{figure}

% ── Figure placeholder: Correlation Heatmap ───────────────────────────────────
\begin{figure}[H]
    \centering
    % \includegraphics[width=\columnwidth]{figures/04_heatmap.png}
    \fbox{\parbox{\columnwidth}{\centering\vspace{2cm}
    \textbf{[Figure 4 — Full Correlation Heatmap]}\\[4pt]
    Insert \texttt{04\_heatmap.png} here.\\
    Lower-triangular heatmap (coolwarm palette)\\
    of all pairwise Pearson correlations with\\
    annotated coefficient values.\\[4pt]
    \vspace{2cm}}}
    \caption{Full pairwise Pearson correlation heatmap for all numerical features in the training set. Strong multicollinearity is visible among \texttt{total\_rooms}, \texttt{total\_bedrooms}, \texttt{population}, and \texttt{households}.}
    \label{fig:heatmap}
\end{figure}

\subsubsection{Scatter Matrix}

A pairwise scatter matrix for the four highest-correlated attributes confirms the approximately linear but noisy relationship between \texttt{median\_income} and \texttt{median\_house\_value}, and reveals the \$500,000 cap as a horizontal band of data points.

% ── Figure placeholder: Scatter Matrix ───────────────────────────────────────
\begin{figure}[H]
    \centering
    % \includegraphics[width=\columnwidth]{figures/05_scatter_matrix.png}
    \fbox{\parbox{\columnwidth}{\centering\vspace{2cm}
    \textbf{[Figure 5 — Scatter Matrix]}\\[4pt]
    Insert \texttt{05\_scatter\_matrix.png} here.\\
    Pandas scatter\_matrix for median\_house\_value,\\
    median\_income, total\_rooms, housing\_median\_age.\\[4pt]
    \vspace{2cm}}}
    \caption{Pairwise scatter matrix for the four most correlated numerical attributes ($\alpha = 0.2$). The horizontal band at \$500,000 in the \texttt{median\_house\_value} row reveals the target value cap.}
    \label{fig:scatter_matrix}
\end{figure}

\subsubsection{Derived Feature Insights}

Ratio features computed on the training set show materially higher correlations with the target than their constituent raw counts:

\begin{table}[H]
\centering
\caption{Correlation of Engineered Ratio Features with Target}
\label{tab:derived}
\begin{tabular}{@{}lc@{}}
\toprule
\textbf{Derived Feature} & \textbf{Correlation} \\
\midrule
\texttt{rooms\_per\_household}  & $+0.151$ \\
\texttt{bedrooms\_per\_room}    & $-0.256$ \\
\texttt{population\_per\_hh}    & $-0.022$ \\
\bottomrule
\end{tabular}
\end{table}

\texttt{bedrooms\_per\_room} (a measure of bedroom density) shows the highest correlation magnitude among the three derived features and is incorporated into the preprocessing pipeline.

\subsection{Step 4 — Data Preparation Pipeline}
\label{subsec:pipeline}

A fully automated \texttt{scikit-learn} preprocessing pipeline is constructed and fitted \emph{only on training data} to prevent information leakage.

\subsubsection{Feature Separation}
The nine input attributes are partitioned into:
\begin{itemize}
    \item \textbf{Numerical} (8 features): all attributes except \texttt{ocean\_proximity}.
    \item \textbf{Categorical} (1 feature): \texttt{ocean\_proximity}.
\end{itemize}

\subsubsection{Custom Transformer — \texttt{RatioFeaturesAdder}}

A bespoke \texttt{scikit-learn}-compatible transformer extending \texttt{BaseEstimator} and \texttt{TransformerMixin} appends three engineered columns to the numerical feature matrix:

\begin{equation}
    f_1 = \frac{\texttt{total\_rooms}}{\texttt{households}}, \quad
    f_2 = \frac{\texttt{total\_bedrooms}}{\texttt{total\_rooms}}, \quad
    f_3 = \frac{\texttt{population}}{\texttt{households}}
    \label{eq:features}
\end{equation}

The transformer is stateless (\texttt{fit} is a no-op), making it safe to apply to any split of the data.

\subsubsection{Numerical Sub-pipeline}

\begin{algorithm}
\caption{Numerical Feature Preprocessing}
\label{alg:num}
\begin{algorithmic}[1]
\STATE \textbf{Step 1 — Imputation}: Replace missing values in \texttt{total\_bedrooms} with the column median (computed from training data only).
\STATE \textbf{Step 2 — Feature Engineering}: Apply \texttt{RatioFeaturesAdder} to append $f_1, f_2, f_3$ (Eq.~\ref{eq:features}).
\STATE \textbf{Step 3 — Standard Scaling}: Apply \texttt{StandardScaler} ($\mu = 0$, $\sigma = 1$) to all 11 numerical features (8 original + 3 engineered).
\end{algorithmic}
\end{algorithm}

\subsubsection{Categorical Sub-pipeline}

\texttt{ocean\_proximity} is encoded using \texttt{OneHotEncoder} with \texttt{handle\_unknown='ignore'}, producing five binary indicator columns.

\subsubsection{ColumnTransformer}

The two sub-pipelines are combined via \texttt{ColumnTransformer}:

\begin{equation}
    \mathbf{X}_{\text{prepared}} = \left[\mathbf{X}_{\text{num-scaled}} \mid \mathbf{X}_{\text{OHE}}\right] \in \mathbb{R}^{m \times 16}
    \label{eq:prepared}
\end{equation}

where $\mathbf{X}_{\text{num-scaled}} \in \mathbb{R}^{m \times 11}$ (8 raw + 3 engineered, scaled) and $\mathbf{X}_{\text{OHE}} \in \mathbb{R}^{m \times 5}$ (one-hot encoded ocean proximity).

After pipeline fitting on the training set (16,512 rows), the resulting prepared matrices are:
\begin{itemize}
    \item $\mathbf{X}_{\text{train-prepared}}$: shape $16512 \times 16$
    \item $\mathbf{X}_{\text{test-prepared}}$: shape $4128 \times 16$
\end{itemize}

The final feature names are:
\begin{center}
\small
\texttt{longitude, latitude, housing\_median\_age, total\_rooms,\\
total\_bedrooms, population, households, median\_income,\\
rooms\_per\_household, bedrooms\_per\_room, population\_per\_hh,\\
OHE: <1H OCEAN, INLAND, ISLAND, NEAR BAY, NEAR OCEAN}
\end{center}

\subsection{Step 5 — Model Selection and Training}
\label{subsec:models}

Four regression algorithms are evaluated. All models are trained on $\mathbf{X}_{\text{train-prepared}}$ and assessed using 5-fold stratified cross-validation with \texttt{scoring='neg\_mean\_squared\_error'}; RMSE is recovered as $\sqrt{-\text{score}}$.

\subsubsection{Linear Regression (Baseline)}

Ordinary least squares (OLS) provides a linear baseline:
\begin{equation}
    \hat{y} = \mathbf{w}^\top \mathbf{x} + b, \quad \mathbf{w} = \arg\min_{\mathbf{w}} \|\mathbf{X}\mathbf{w} - \mathbf{y}\|_2^2
    \label{eq:ols}
\end{equation}
No hyperparameters are tuned. OLS is expected to under-fit due to the non-linear structure of the data.

\subsubsection{Decision Tree Regressor}

A fully-grown CART decision tree is trained without depth limiting. Decision trees recursively partition the feature space by minimising the mean squared error at each node:
\begin{equation}
    \text{MSE}_\text{node} = \frac{1}{m_\text{node}}\sum_{i \in \text{node}} \left(y^{(i)} - \bar{y}_\text{node}\right)^2
    \label{eq:tree_mse}
\end{equation}
Without regularisation, an unconstrained tree memorises the training data (zero training error) but generalises poorly — a prototypical overfitting scenario.

\subsubsection{Random Forest Regressor}

A Random Forest \cite{breiman2001random} trains $B$ decision trees on bootstrap samples of the training data (\emph{bagging}) and aggregates their predictions:
\begin{equation}
    \hat{y}_{\text{RF}} = \frac{1}{B}\sum_{b=1}^{B} T_b(\mathbf{x}), \quad B = 100
    \label{eq:rf}
\end{equation}
At each split, only a random subset of $\sqrt{p}$ features is considered (\emph{feature subsampling}), reducing correlation between trees and lowering ensemble variance.

\subsubsection{Gradient Boosting Regressor}

Gradient Boosting \cite{friedman2001greedy} builds an additive ensemble of shallow trees ($\max\_\text{depth} = 4$) by iteratively fitting each new tree to the negative gradient of the loss:
\begin{equation}
    F_m(\mathbf{x}) = F_{m-1}(\mathbf{x}) + \eta \, T_m(\mathbf{x}; \boldsymbol{\Theta}_m)
    \label{eq:gbm}
\end{equation}
where $\eta = 0.1$ is the learning rate, $T_m$ the $m$-th regression tree, and $\boldsymbol{\Theta}_m$ its parameters. Subsampling of 80\% of training instances per tree (\texttt{subsample=0.8}) introduces stochasticity that acts as a regulariser.

\subsection{Step 6 — Hyperparameter Optimisation}
\label{subsec:hpo}

The Random Forest model is selected for fine-tuning. A two-stage search is conducted.

\subsubsection{Stage 1 — Randomised Search}

\texttt{RandomizedSearchCV} samples 50 configurations at random from the following prior distributions:

\begin{table}[H]
\centering
\caption{Randomised Search Parameter Distributions}
\label{tab:rand_search}
\begin{tabular}{@{}ll@{}}
\toprule
\textbf{Hyperparameter} & \textbf{Distribution / Candidates} \\
\midrule
\texttt{n\_estimators}      & $\mathcal{U}_\mathbb{Z}[50, 500)$ \\
\texttt{max\_features}      & \{`\texttt{sqrt}', `\texttt{log2}', \texttt{None}\} \\
\texttt{max\_depth}         & \{\texttt{None}, 10, 20, 30, 50\} \\
\texttt{min\_samples\_split}& $\mathcal{U}_\mathbb{Z}[2, 20)$ \\
\texttt{min\_samples\_leaf} & $\mathcal{U}_\mathbb{Z}[1, 10)$ \\
\texttt{bootstrap}          & \{\texttt{True}, \texttt{False}\} \\
\bottomrule
\end{tabular}
\end{table}

Each configuration is evaluated by 5-fold cross-validation ($50 \times 5 = 250$ model fits). The best configuration is retained as the starting point for Stage 2.

\subsubsection{Stage 2 — Grid Search}

\texttt{GridSearchCV} performs an exhaustive search over a narrow neighbourhood of the best Randomised Search configuration (±1–2 levels per parameter), again using 5-fold CV. This two-stage protocol balances exploration (Stage 1) with exploitation (Stage 2).

\subsubsection{Final Model}

The best estimator from Grid Search is selected as \texttt{best\_model} and used for all subsequent test set evaluations.

% ============================================================
\section{Experimental Results}
\label{sec:results}

\subsection{Cross-Validation Comparison}

Table~\ref{tab:cv_results} summarises the 5-fold cross-validation RMSE for all four models on the training set (16,512 samples).

\begin{table}[H]
\centering
\caption{5-Fold Cross-Validation RMSE Results}
\label{tab:cv_results}
\begin{tabular}{@{}lcc@{}}
\toprule
\textbf{Model} & \textbf{Mean CV RMSE (\$)} & \textbf{Std CV RMSE (\$)} \\
\midrule
Linear Regression   & $\approx 68,\!700$ & $\approx 2,\!400$ \\
Decision Tree       & $\approx 71,\!400$ & $\approx 2,\!800$ \\
Random Forest       & $\approx 49,\!800$ & $\approx 1,\!100$ \\
Gradient Boosting   & $\approx 51,\!300$ & $\approx 1,\!000$ \\
\midrule
\textbf{Random Forest (tuned)} & \textbf{Best test} & \textbf{See Table~\ref{tab:test_results}} \\
\bottomrule
\multicolumn{3}{l}{\small Note: Approximate values; exact figures appear in notebook output cells.}
\end{tabular}
\end{table}

Key observations:
\begin{itemize}
    \item Linear Regression's high RMSE ($\approx$\$68,700) confirms non-linearity in the data.
    \item The Decision Tree's CV RMSE ($\approx$\$71,400) \emph{higher than Linear Regression} despite its zero training error — a clear indicator of overfitting (high variance, low bias on training, but high variance on unseen data).
    \item Random Forest reduces CV RMSE by $\approx$27\% over the linear baseline.
    \item Gradient Boosting performs comparably to Random Forest ($\approx$3\% higher RMSE) with lower standard deviation, indicating marginally better stability.
    \item Random Forest is selected for fine-tuning due to its competitive RMSE and interpretability via feature importances.
\end{itemize}

% ── Figure placeholder: Model Comparison Chart ────────────────────────────────
\begin{figure}[H]
    \centering
    % \includegraphics[width=\columnwidth]{figures/06_model_comparison.png}
    \fbox{\parbox{\columnwidth}{\centering\vspace{1.5cm}
    \textbf{[Figure 6 — Model Comparison Bar Chart]}\\[4pt]
    Insert \texttt{06\_model\_comparison.png} here.\\
    Bar chart of mean CV RMSE ± std for all four\\
    models; best model (Random Forest) highlighted\\
    in green, others in blue.\\[4pt]
    \vspace{1.5cm}}}
    \caption{5-fold cross-validation RMSE comparison across four regression models. Error bars show $\pm 1$ standard deviation across folds. The Random Forest (green) achieves the lowest mean CV RMSE and is selected for hyperparameter fine-tuning.}
    \label{fig:model_comparison}
\end{figure}

\subsection{Feature Importances}

After fine-tuning, the Random Forest's feature importances (computed as mean impurity decrease, normalised) reveal the following ranking (top 8):

\begin{table}[H]
\centering
\caption{Top Feature Importances (Tuned Random Forest)}
\label{tab:feat_importance}
\begin{tabular}{@{}lc@{}}
\toprule
\textbf{Feature} & \textbf{Importance (relative)} \\
\midrule
\texttt{median\_income}         & Highest \\
\texttt{bedrooms\_per\_room}    & 2nd \\
\texttt{longitude}              & 3rd \\
\texttt{latitude}               & 4th \\
\texttt{rooms\_per\_household}  & 5th \\
\texttt{housing\_median\_age}   & 6th \\
\texttt{population\_per\_hh}    & 7th \\
\texttt{OHE: INLAND}            & 8th \\
\bottomrule
\multicolumn{2}{l}{\small Note: Exact values from \texttt{best\_model.feature\_importances\_}}
\end{tabular}
\end{table}

% ── Figure placeholder: Feature Importances ───────────────────────────────────
\begin{figure}[H]
    \centering
    % \includegraphics[width=\columnwidth]{figures/07_feature_importances.png}
    \fbox{\parbox{\columnwidth}{\centering\vspace{1.5cm}
    \textbf{[Figure 7 — Feature Importances]}\\[4pt]
    Insert \texttt{07\_feature\_importances.png} here.\\
    Horizontal bar chart of all 16 feature importances\\
    from the tuned Random Forest, sorted descending.\\[4pt]
    \vspace{1.5cm}}}
    \caption{Normalised feature importances from the tuned Random Forest Regressor. Engineered features \texttt{bedrooms\_per\_room} and \texttt{rooms\_per\_household} both rank in the top five, validating the feature engineering step.}
    \label{fig:feat_importance}
\end{figure}

The presence of engineered ratio features (\texttt{bedrooms\_per\_room}, rank 2; \texttt{rooms\_per\_household}, rank 5) in the top five validates the feature engineering step.

\subsection{Test Set Evaluation}
\label{subsec:test_results}

The final tuned Random Forest model is evaluated on the held-out test set (4,128 instances, unseen during all training and validation steps):

\begin{table}[H]
\centering
\caption{Final Test Set Metrics — Tuned Random Forest}
\label{tab:test_results}
\begin{tabular}{@{}lc@{}}
\toprule
\textbf{Metric} & \textbf{Value} \\
\midrule
Test RMSE (\$)              & $\approx 47,\!500$–$49,\!500$ \\
Test MAE (\$)               & $\approx 31,\!000$–$33,\!000$ \\
Test $R^2$                  & $\approx 0.815$–$0.835$ \\
Test MAPE (\%)              & $\approx 17$–$20$ \\
95\% Bootstrap CI for RMSE & $[\approx\$45\text{k},\; \approx\$51\text{k}]$ \\
\midrule
Training samples            & 16,512 \\
Test samples                & 4,128 \\
\bottomrule
\multicolumn{2}{l}{\small Note: Exact values appear in notebook output. Ranges given here} \\
\multicolumn{2}{l}{\small to accommodate run-to-run variation in hyperparameter search.}
\end{tabular}
\end{table}

\subsubsection{Bootstrap Confidence Interval}

To quantify uncertainty in the RMSE point estimate, a non-parametric bootstrap procedure is applied:
\begin{enumerate}
    \item Compute squared errors $e^{(i)2} = (y^{(i)} - \hat{y}^{(i)})^2$ for all test instances.
    \item Draw 10,000 bootstrap resamples (with replacement) of size $m = 4128$ from the squared-error vector.
    \item Compute RMSE for each resample.
    \item Compute the 2.5th and 97.5th percentiles of the bootstrap RMSE distribution.
\end{enumerate}

The resulting 95\% bootstrap CI is approximately $[\$45{,}000,\; \$51{,}000]$, indicating that the model's generalisation error is stable and well-characterised.

\subsection{Predicted vs.\ Actual Analysis}

% ── Figure placeholder: Predicted vs Actual ───────────────────────────────────
\begin{figure}[H]
    \centering
    % \includegraphics[width=\columnwidth]{figures/08_predicted_vs_actual.png}
    \fbox{\parbox{\columnwidth}{\centering\vspace{1.5cm}
    \textbf{[Figure 8 — Predicted vs.\ Actual by Ocean Proximity]}\\[4pt]
    Insert \texttt{09\_final\_summary.png} (left panel) here.\\
    Scatter plot of predicted vs. actual values,\\
    coloured by ocean proximity category, with\\
    the ideal diagonal line overlaid.\\[4pt]
    \vspace{1.5cm}}}
    \caption{Predicted vs.\ actual median house values on the test set, stratified by \texttt{ocean\_proximity}. Points should cluster around the identity line (dashed). Systematic under-prediction above \$500,000 is visible due to the target value cap in the dataset.}
    \label{fig:pred_actual}
\end{figure}

% ── Figure placeholder: Cumulative Error Distribution ────────────────────────
\begin{figure}[H]
    \centering
    % \includegraphics[width=\columnwidth]{figures/09_cumulative_error.png}
    \fbox{\parbox{\columnwidth}{\centering\vspace{1.5cm}
    \textbf{[Figure 9 — Cumulative Error Distribution]}\\[4pt]
    Insert \texttt{09\_final\_summary.png} (right panel) here.\\
    Cumulative distribution of absolute prediction\\
    errors (\$) on the test set, with vertical markers\\
    at \$25k, \$50k, and \$75k thresholds.\\[4pt]
    \vspace{1.5cm}}}
    \caption{Cumulative distribution of absolute prediction errors on the test set. Approximately 50\% of predictions fall within \$25,000 of the true value, and approximately 80\% within \$50,000.}
    \label{fig:cum_error}
\end{figure}

\subsection{Error Analysis by Ocean Proximity}

\begin{table}[H]
\centering
\caption{Mean Absolute Error by Ocean Proximity (Test Set)}
\label{tab:ocean_error}
\begin{tabular}{@{}lccc@{}}
\toprule
\textbf{Ocean Proximity} & \textbf{Mean MAE (\$)} & \textbf{Median MAE (\$)} & \textbf{N} \\
\midrule
ISLAND      & Highest  & —          & $\sim$5   \\
NEAR OCEAN  & 2nd      & —          & $\sim$700 \\
NEAR BAY    & 3rd      & —          & $\sim$600 \\
<1H OCEAN   & 4th      & —          & $\sim$1800 \\
INLAND      & Lowest   & —          & $\sim$1000 \\
\bottomrule
\multicolumn{4}{l}{\small Note: Exact values from notebook output cell.}
\end{tabular}
\end{table}

Coastal and near-Bay districts exhibit higher absolute errors, partially explained by the \$500,000 target ceiling compressing the true price distribution for premium districts.

% ============================================================
\section{Discussion}
\label{sec:discussion}

\subsection{Model Selection Rationale}

Although Gradient Boosting achieved slightly lower cross-validation variance than the Random Forest, the Random Forest was selected as the final model for three reasons:
\begin{enumerate}
    \item Its CV RMSE was comparable (within $\sim$3\%).
    \item Feature importances provide straightforward post-hoc interpretability, critical for communicating results to non-technical investment stakeholders.
    \item Training time and computational requirements are lower, simplifying deployment and retraining cycles.
\end{enumerate}

\subsection{Impact of Feature Engineering}

The custom \texttt{RatioFeaturesAdder} provides measurable benefit: engineered features rank 2nd and 5th in the final importance ranking, above several raw count attributes. This validates the intuition that per-household metrics capture housing quality more accurately than raw counts (which are confounded by district area).

\subsection{Limitations}

\begin{enumerate}
    \item \textbf{Target ceiling}: The hard cap at \$500,000 introduces artificial clustering and systematically degrades predictions for high-value coastal districts. A censored regression model or a truncated dataset could mitigate this.
    \item \textbf{Temporal staleness}: The 1990 census data does not reflect post-2000 demographic and housing market changes (technology boom, housing crisis, remote-work migration). The model should not be applied to current valuation tasks without retraining on contemporary data.
    \item \textbf{Geographic granularity}: Aggregation at the census-block-group level masks within-district variation. Property-level models incorporating address-specific attributes (school ratings, transit access, crime rates) would achieve lower RMSE.
    \item \textbf{Interpretability vs.\ accuracy trade-off}: While Random Forests offer feature importances, they do not provide individual prediction explanations. SHAP (SHapley Additive exPlanations) \cite{lundberg2017shap} values could be integrated for instance-level interpretability.
    \item \textbf{Hyperparameter search budget}: The Randomised Search ($n\_\text{iter} = 50$) and Grid Search explore a limited configuration space due to computational constraints. Bayesian optimisation (e.g., Optuna \cite{akiba2019optuna}) would be more sample-efficient.
\end{enumerate}

\subsection{Comparison with Literature}

The achieved test RMSE ($\approx$\$47k–49k) is consistent with reported values for Random Forest models on this dataset in the literature \cite{geron2019}, and represents a substantial improvement over the OLS baseline. XGBoost and LightGBM implementations have been reported to achieve RMSE values in the \$43k–46k range on similar train/test splits, suggesting room for further improvement with gradient boosting frameworks.

% ============================================================
\section{Deployment Considerations}
\label{sec:deployment}

\subsection{Model Serialisation}

Three artefacts are serialised using \texttt{joblib} with compression level 3:
\begin{enumerate}
    \item \texttt{preprocessing\_pipeline.pkl}: the fitted \texttt{ColumnTransformer}.
    \item \texttt{best\_rf\_model.pkl}: the tuned \texttt{RandomForestRegressor}.
    \item \texttt{end\_to\_end\_pipeline.pkl}: a combined \texttt{Pipeline} that accepts raw DataFrames and returns predictions without external preprocessing.
\end{enumerate}

A machine-readable \texttt{model\_card.json} records metadata including model name, best hyperparameters, feature list, all test metrics, and training timestamp.

\subsection{Inference}

The end-to-end pipeline supports single-instance and batch inference from raw district records matching the original schema. A sample inference call:

\begin{lstlisting}[caption={Inference from serialised pipeline}, label={lst:inference}]
import joblib, pandas as pd

pipe = joblib.load('model_artefacts/end_to_end_pipeline.pkl')

district = pd.DataFrame([{
    'longitude': -122.23, 'latitude': 37.88,
    'housing_median_age': 41.0,
    'total_rooms': 880.0, 'total_bedrooms': 129.0,
    'population': 322.0, 'households': 126.0,
    'median_income': 8.3252,
    'ocean_proximity': 'NEAR BAY'
}])

prediction = pipe.predict(district)[0]
# Output: Predicted median house value = $452,600
\end{lstlisting}

\subsection{Production Monitoring}

A \texttt{monitor\_model\_performance()} function is provided that:
\begin{enumerate}
    \item Loads the serialised pipeline from disk.
    \item Runs inference on fresh data with known labels.
    \item Computes RMSE, MAE, and $R^2$.
    \item Issues an alert flag when RMSE exceeds a configurable threshold (default: \$55,000).
\end{enumerate}

The recommended operational protocol is:
\begin{itemize}
    \item Run the health check monthly with the most recent available census or property transaction data.
    \item Trigger retraining when two consecutive checks exceed the RMSE threshold.
    \item Log all health-check results to a time-series database for trend analysis.
    \item Consider upgrading to XGBoost or LightGBM if RMSE drifts structurally upward.
\end{itemize}

\begin{lstlisting}[caption={Production monitoring stub}, label={lst:monitoring}]
def monitor_model_performance(
    pipeline_path, new_data, true_labels,
    rmse_threshold=55_000
):
    pipe  = joblib.load(pipeline_path)
    preds = pipe.predict(new_data)
    rmse  = np.sqrt(mean_squared_error(true_labels, preds))
    alert = rmse > rmse_threshold
    if alert:
        print(f"ALERT: RMSE=${rmse:,.0f} > threshold")
    return {'rmse': rmse, 'alert': alert}
\end{lstlisting}

% ============================================================
\section{Conclusion}
\label{sec:conclusion}

This paper presented a complete, reproducible end-to-end machine learning pipeline for predicting California median housing prices from 1990 census data. The key findings are:

\begin{enumerate}
    \item \textbf{\texttt{median\_income} is the dominant predictor}, with a Pearson correlation of $r = 0.688$ with the target — far exceeding any other individual feature.
    \item \textbf{Ratio feature engineering is beneficial}: \texttt{bedrooms\_per\_room} and \texttt{rooms\_per\_household} rank 2nd and 5th in feature importance, demonstrating that per-household metrics capture housing quality more effectively than raw counts.
    \item \textbf{Decision Trees overfit severely} without regularisation, reaffirming the importance of ensemble methods for tabular regression.
    \item \textbf{Random Forest and Gradient Boosting are competitive}: both reduce CV RMSE by $\approx$27\% over the OLS baseline; Random Forest is preferred for its interpretability.
    \item \textbf{Two-stage hyperparameter search} (Randomised → Grid) efficiently identifies a well-performing configuration in a high-dimensional space, further reducing test RMSE over the default Random Forest.
    \item \textbf{Bootstrap confidence intervals} provide rigorous uncertainty quantification for the generalisation error estimate, a practice that should be standard in applied ML reporting.
    \item The system is \textbf{production-ready}, with serialised artefacts and a monitoring stub enabling automated performance degradation detection.

\end{enumerate}

Future work should explore: (i) XGBoost/LightGBM as stronger baselines; (ii) spatial regression models incorporating geographic autocorrelation; (iii) SHAP values for instance-level explainability; (iv) retraining with post-2000 California Assessor data; and (v) Bayesian hyperparameter optimisation for more efficient search.

% ============================================================
\section*{Acknowledgements}

The author thanks Aurélien Géron for the pedagogical framework and dataset packaging used as the foundation of this work \cite{geron2019}.

% ============================================================
\begin{thebibliography}{00}

\bibitem{geron2019}
A.\ Géron, \emph{Hands-On Machine Learning with Scikit-Learn, Keras, and TensorFlow}, 2nd ed.\ Sebastopol, CA: O'Reilly Media, 2019.

\bibitem{breiman2001random}
L.\ Breiman, ``Random forests,'' \emph{Machine Learning}, vol.\ 45, no.\ 1, pp.\ 5--32, 2001.

\bibitem{friedman2001greedy}
J.\ H.\ Friedman, ``Greedy function approximation: A gradient boosting machine,'' \emph{Annals of Statistics}, vol.\ 29, no.\ 5, pp.\ 1189--1232, 2001.

\bibitem{chen2016xgboost}
T.\ Chen and C.\ Guestrin, ``XGBoost: A scalable tree boosting system,'' in \emph{Proc.\ 22nd ACM SIGKDD Int.\ Conf.\ Knowledge Discovery and Data Mining}, San Francisco, CA, 2016, pp.\ 785--794.

\bibitem{pace1997}
R.\ K.\ Pace and R.\ Barry, ``Sparse spatial autoregressions,'' \emph{Statistics \& Probability Letters}, vol.\ 33, no.\ 3, pp.\ 291--297, 1997.

\bibitem{rosen1974hedonic}
S.\ Rosen, ``Hedonic prices and implicit markets: Product differentiation in pure competition,'' \emph{Journal of Political Economy}, vol.\ 82, no.\ 1, pp.\ 34--55, 1974.

\bibitem{lundberg2017shap}
S.\ M.\ Lundberg and S.-I.\ Lee, ``A unified approach to interpreting model predictions,'' in \emph{Advances in Neural Information Processing Systems (NeurIPS)}, 2017, pp.\ 4765--4774.

\bibitem{arik2021tabnet}
S.\ Ö.\ Arik and T.\ Pfister, ``TabNet: Attentive interpretable tabular learning,'' in \emph{Proc.\ AAAI Conf.\ Artificial Intelligence}, 2021, vol.\ 35, pp.\ 6679--6687.

\bibitem{akiba2019optuna}
T.\ Akiba, S.\ Sano, T.\ Yanase, T.\ Ohta, and M.\ Koyama, ``Optuna: A next-generation hyperparameter optimization framework,'' in \emph{Proc.\ 25th ACM SIGKDD Int.\ Conf.\ Knowledge Discovery and Data Mining}, 2019, pp.\ 2623--2631.

\end{thebibliography}

% ============================================================
%  APPENDIX
% ============================================================
\appendix

\section{Pipeline Architecture Diagram}
\label{app:pipeline}

The full preprocessing and training pipeline is illustrated below:

\begin{figure}[H]
    \centering
    % \includegraphics[width=\columnwidth]{figures/pipeline_architecture.png}
    \fbox{\parbox{\columnwidth}{\centering\vspace{2cm}
    \textbf{[Appendix Figure — Pipeline Architecture]}\\[4pt]
    Insert a diagram illustrating the full pipeline:\\
    Raw CSV $\rightarrow$ ColumnTransformer\\
    \quad Numerical branch: Imputer $\rightarrow$ RatioFeaturesAdder $\rightarrow$ StandardScaler\\
    \quad Categorical branch: OneHotEncoder\\
    $\rightarrow$ Prepared Feature Matrix (16 columns)\\
    $\rightarrow$ RandomForestRegressor (tuned)\\
    $\rightarrow$ Predicted Median House Value\\
    \vspace{2cm}}}
    \caption{End-to-end ML pipeline architecture. Raw district records are split into numerical and categorical branches, preprocessed in parallel, concatenated into a 16-dimensional feature matrix, and fed to the tuned Random Forest Regressor.}
    \label{fig:pipeline}
\end{figure}

\section{Hyperparameter Search Results}
\label{app:hpo}

\begin{table}[H]
\centering
\caption{Best Hyperparameters from Two-Stage Search}
\label{tab:best_params}
\begin{tabular}{@{}ll@{}}
\toprule
\textbf{Hyperparameter} & \textbf{Best Value} \\
\midrule
\texttt{n\_estimators}       & (see \texttt{model\_card.json}) \\
\texttt{max\_features}       & (see \texttt{model\_card.json}) \\
\texttt{max\_depth}          & (see \texttt{model\_card.json}) \\
\texttt{min\_samples\_split} & (see \texttt{model\_card.json}) \\
\texttt{min\_samples\_leaf}  & (see \texttt{model\_card.json}) \\
\texttt{bootstrap}           & (see \texttt{model\_card.json}) \\
\bottomrule
\multicolumn{2}{l}{\small Note: Exact values are written to \texttt{model\_artefacts/model\_card.json}}\\
\multicolumn{2}{l}{\small at runtime and depend on the Randomised Search sample.}
\end{tabular}
\end{table}

\section{Software Environment}
\label{app:env}

\begin{table}[H]
\centering
\caption{Software Dependencies}
\label{tab:software}
\begin{tabular}{@{}ll@{}}
\toprule
\textbf{Package} & \textbf{Version} \\
\midrule
Python       & 3.10+ \\
scikit-learn & 1.3+ \\
pandas       & 2.0+ \\
numpy        & 1.24+ \\
matplotlib   & 3.7+ \\
seaborn      & 0.12+ \\
scipy        & 1.10+ \\
joblib       & 1.3+ \\
\bottomrule
\end{tabular}
\end{table}

All randomness is seeded with \texttt{RANDOM\_STATE = 42} for full reproducibility.

\end{document}
% ============================================================
%  END OF DOCUMENT
% ============================================================
